\documentclass[a4paper]{article}

\usepackage[fontset=fandol]{ctex}
\usepackage{amsmath, amssymb}
\usepackage{graphicx}
\usepackage[margin=2.45cm]{geometry}
\usepackage[most]{tcolorbox}
\usepackage{hyperref}
\usepackage{booktabs}
\usepackage{float}
\usepackage{array}
\usepackage{xcolor}
\usepackage{enumitem}
\usepackage{caption}

\setlength{\parindent}{2em}
\setlength{\parskip}{0.35em}
\setlength{\emergencystretch}{3em}
\linespread{1.06}
\sloppy
\raggedbottom
\vfuzz=4pt

\newcolumntype{L}[1]{>{\raggedright\arraybackslash}p{#1}}
\newcolumntype{C}[1]{>{\centering\arraybackslash}p{#1}}

\DeclareMathOperator*{\argmin}{arg\,min}
\DeclareMathOperator*{\argmax}{arg\,max}
\DeclareMathOperator{\CE}{CE}
\newcommand{\E}{\mathbb{E}}
\newcommand{\st}{s}
\newcommand{\at}{a}
\newcommand{\traj}{\tau}
\newcommand{\policy}{\pi_\theta}

\hypersetup{
    colorlinks=true,
    linkcolor=blue!60!black,
    urlcolor=blue!60!black
}

\setlist[itemize]{leftmargin=2.2em,itemsep=0.22em,topsep=0.25em}
\setlist[enumerate]{leftmargin=2.4em,itemsep=0.22em,topsep=0.25em}
\setlist[description]{leftmargin=3.0em,itemsep=0.18em,topsep=0.25em}

\newtcolorbox{knowledgebox}[1]{
    enhanced, colback=blue!5!white, colframe=blue!75!black, colbacktitle=blue!75!black,
    coltitle=white, fonttitle=\bfseries, title={#1},
    attach boxed title to top left={yshift=-2mm, xshift=2mm},
    boxrule=1pt, sharp corners
}

\newtcolorbox{importantbox}[1]{
    enhanced, colback=yellow!10!white, colframe=yellow!80!black, colbacktitle=yellow!80!black,
    coltitle=black, fonttitle=\bfseries, title={#1}, sharp corners
}

\newtcolorbox{warningbox}[1]{
    enhanced, colback=red!5!white, colframe=red!75!black, colbacktitle=red!75!black,
    coltitle=white, fonttitle=\bfseries, title={#1}, sharp corners
}

\newcommand{\notetitle}{强化学习概述（一）：RL 也是机器学习三步骤}
\newcommand{\notesubtitle}{Actor, Environment, Reward, Return, and the Intuition of Policy Gradient}
\newcommand{\noteauthors}{基于李宏毅老师《机器学习 2025》课程内容整理}
\newcommand{\notedate}{\today}
\newcommand{\videochannel}{李宏毅深度学习课堂（Bilibili）}
\newcommand{\videoduration}{约 45 分 01 秒}
\newcommand{\videourl}{https://www.bilibili.com/video/BV1TAtwzTE1S?p=46}

\newcommand{\framefig}[4]{%
\begin{figure}[H]
    \centering
    \includegraphics[width=#1\textwidth]{#2}
    \caption{#3\protect\footnotemark}
\end{figure}
\footnotetext{视频画面时间区间：#4。}
}

\begin{document}
\pagestyle{empty}

\begin{center}
    \vspace*{1.0cm}
    {\Huge\bfseries \notetitle\par}
    \vspace{0.35cm}
    {\LARGE \notesubtitle\par}
    \vspace{0.75cm}
    \includegraphics[width=0.62\textwidth]{video.jpg}\par
    \vspace{0.75cm}
    {\large \noteauthors\par}
    {\large \notedate\par}
    \vspace{0.75cm}
    \begin{tcolorbox}[width=0.86\textwidth,center,sharp corners,
                      colback=black!3!white,colframe=black!50!white]
        \centering
        \textbf{视频信息}\\[3pt]
        频道：\videochannel \\
        时长：\videoduration \\
        链接：\href{\videourl}{\videourl}
    \end{tcolorbox}
\end{center}

\newpage
\tableofcontents
\newpage
\pagestyle{plain}

% =====================================================================
\section{从监督学习走向强化学习}
% =====================================================================

这节课开始进入 deep reinforcement learning。老师一开始刻意把强化学习拉回前面熟悉的框架：机器学习并不是一堆完全不同的魔法，而是常常可以被理解成“找一个函数”。在监督学习中，我们有输入、模型输出和人类标注；模型的目标是让输出接近 label。强化学习也在找函数，只是 label 不再是每一步直接给出的正确答案，而是来自模型与环境互动之后得到的 reward。

\framefig{0.92}{figures/fig01_supervised_to_rl_label_difficulty.jpg}
{从监督学习到 RL：猫狗分类可以有人类 label，但围棋某个盘面的最佳下一手很难直接标注。}
{00:03:30--00:03:45}

围棋例子说明了 RL 的动机。我们可以收集高手棋谱，把高手某一步当成 label；但那一步是否真的是该局面下最好的动作，未必可知。更一般地，在许多任务中，人类很难告诉机器“此刻正确动作是什么”，却可以让机器行动以后判断结果好坏。例如游戏赢了还是输了、得了多少分、机器人是否完成任务。这种“没有逐步标签，但有互动反馈”的场景，就是强化学习适合登场的地方。

\framefig{0.92}{figures/fig02_outline_rl_topics.jpg}
{本节强化学习大纲：先解释 RL 与机器学习三步骤的对应，再进入 policy gradient、actor-critic、reward shaping 等主题。}
{00:04:15--00:04:30}

\begin{importantbox}{RL 与监督学习的关键差别}
    监督学习通常给定 $(x,y)$，模型只要让 $f_\theta(x)$ 接近 $y$。强化学习中，模型先根据观察采取动作，动作会改变环境，环境再回传新的观察和 reward。模型没有每一步的标准答案，只有长期互动后的好坏信号。
\end{importantbox}

\subsection{RL 中的四个基本对象}

强化学习的基本循环包含 actor、environment、observation 和 reward。actor 是我们要学习的函数，也常被称为 policy。environment 是 actor 互动的外部世界。environment 给 actor 一个 observation，actor 根据 observation 采取 action；action 影响 environment，environment 再给出新的 observation 和 reward。

\framefig{0.92}{figures/fig03_actor_environment_loop.jpg}
{RL 的基本循环：environment 给 observation，actor 输出 action，action 影响环境，环境再回传 reward 与新 observation。}
{00:07:00--00:07:15}

若用符号表示，policy 是一个从状态或观察到动作分布的函数：
\[
\policy(a\mid s)=P_\theta(a_t=a\mid s_t=s).
\]
其中 $s_t$ 是第 $t$ 步的 observation 或 state，$a_t$ 是 actor 采取的 action，$\theta$ 是 policy network 的参数。若动作空间是离散的，$\policy(a\mid s)$ 可以理解为神经网络输出的 softmax 概率。

\begin{knowledgebox}{observation 与 state}
    课程中主要用 observation 解释，因为游戏画面或棋盘局面就是 actor 实际看到的输入。严格 RL 文献中，state 往往指环境内部的完整状态，observation 指 agent 可见的信息。若环境完全可观测，二者可以近似等同；若不完全可观测，observation 只是 state 的一部分投影。
\end{knowledgebox}

\subsection{本章小结}

强化学习仍然可以被看成“找一个函数”，但这个函数不是把输入直接映射到固定 label，而是在环境中持续行动，并通过 reward 知道行动好坏。RL 的难点来自 label 的缺失、动作对未来观察的影响，以及 reward 可能很延迟。

% =====================================================================
\section{例子：游戏、太空侵略者与 AlphaGo}
% =====================================================================

为了把抽象循环具体化，课程先用 Atari 游戏 Space Invader 作例子。actor 看到游戏画面，动作可以是左移、右移或开火。游戏画面中的外星人、防护罩、飞船位置、分数等都组成 observation。actor 的目标不是预测某个固定 label，而是通过一连串动作获得尽可能高的分数。

\framefig{0.92}{figures/fig04_space_invader_state_action.jpg}
{Space Invader 中，actor 的观察是游戏画面，动作包括左移、右移与开火；目标是杀掉外星人并保护飞船。}
{00:08:00--00:08:15}

当 actor 开火并击中外星人，environment 会给 reward，例如 $r=5$。若动作没有立刻带来分数，reward 可能是 $0$。游戏结束时如果飞船被摧毁，也可能得到负面结果。重点是：actor 必须学习的不只是“当前哪一个动作看起来合理”，而是“哪些动作序列能带来更高的总 reward”。

\framefig{0.92}{figures/fig05_space_invader_reward_objective.jpg}
{在游戏中，actor 根据画面采取动作，若击中外星人就得到 reward；目标是最大化整局游戏的期望 reward。}
{00:11:00--00:11:15}

围棋的例子则强调 delayed reward。AlphaGo 的 observation 是棋盘上的黑白子位置，action 是下一步落子位置。environment 可以看成对手与围棋规则。大多数中间步骤没有即时 reward，只有整局结束时，赢了得到正 reward，输了得到负 reward。

\framefig{0.92}{figures/fig06_alphago_observation_action.jpg}
{AlphaGo 设定中，observation 是棋盘局面，action 是下一步落子，environment 包含对手与规则。}
{00:12:15--00:12:30}

\framefig{0.92}{figures/fig07_go_sparse_terminal_reward.jpg}
{围棋的 reward 很稀疏：多数中间动作 reward 为 0，只有游戏结束时根据胜负给出最终 reward。}
{00:14:00--00:14:15}

\begin{warningbox}{delayed reward 带来的 credit assignment}
    若最后赢棋，究竟是哪几步促成胜利？若最后输棋，哪一步应该负责？强化学习必须处理 credit assignment：把长期结果分配回早先动作。围棋这类稀疏 reward 任务尤其困难，因为中间几乎没有逐步指导。
\end{warningbox}

\subsection{本章小结}

Space Invader 展示了比较直观的即时 reward：击中外星人就加分。围棋展示了更困难的稀疏与延迟 reward：只有整局结束时才知道胜负。两者共同说明，RL 关心的是一连串动作造成的长期结果，而不是单一步骤的局部正确性。

% =====================================================================
\section{RL 也是机器学习三步骤}
% =====================================================================

老师接着把 RL 放回机器学习三步骤：第一步，写出带未知参数的函数；第二步，定义 loss 或目标函数；第三步，用优化方法调整参数。强化学习看起来复杂，但仍然可以被整理成这三个步骤。

\framefig{0.92}{figures/fig08_three_steps_machine_learning.jpg}
{机器学习三步骤同样适用于 RL：定义含未知参数的函数、定义目标，再进行优化。}
{00:14:15--00:14:30}

\subsection{Step 1：定义 actor，也就是 policy network}

第一步是定义 policy network。以游戏为例，输入可以是画面像素矩阵，输出层每个神经元对应一个可选动作，例如 left、right、fire。网络输出不一定直接选择最大分数动作，而可以形成一个动作分布，再从中采样。

\framefig{0.92}{figures/fig09_policy_network_scores.jpg}
{policy network 将 observation 映射到动作分数或概率；每个输出神经元对应一个可执行 action。}
{00:16:00--00:16:15}

若输出是 logits $z_\theta(s,a)$，可以通过 softmax 得到 policy：
\[
\pi_\theta(a\mid s)
=\frac{\exp(z_\theta(s,a))}
{\sum_{a'\in\mathcal{A}}\exp(z_\theta(s,a'))}.
\]
其中 $\mathcal{A}$ 是动作集合。实际执行时，可以选择
\[
a_t\sim \pi_\theta(\cdot\mid s_t),
\]
也就是按概率采样动作，而不总是取 $\argmax_a \pi_\theta(a\mid s_t)$。

\begin{knowledgebox}{为什么要采样动作}
    采样让 policy 保持探索能力。若一开始总是选当前分数最高的动作，模型可能永远没有机会发现其他动作其实更好。强化学习中的 exploration 与 exploitation 平衡，正是后续算法设计的重要问题。
\end{knowledgebox}

\subsection{Step 2：定义目标，最大化 return}

第二步是定义要最大化的量。与监督学习不同，RL 不是直接给每个 state-action pair 一个正确 label，而是观察一整段互动。一次从开始到结束的互动称为 episode；在 episode 中得到的 reward 总和称为 return。

\framefig{0.92}{figures/fig11_episode_return.jpg}
{一整局游戏形成一个 episode；从开始到结束的 reward 总和称为 return。}
{00:22:15--00:22:30}

设一条 trajectory 为
\[
\tau=(s_1,a_1,r_1,s_2,a_2,r_2,\ldots,s_T,a_T,r_T).
\]
不考虑折扣时，return 是
\[
R(\tau)=\sum_{t=1}^{T}r_t.
\]
很多 RL 文献也使用折扣回报：
\[
R_\gamma(\tau)=\sum_{t=1}^{T}\gamma^{t-1}r_t,
\qquad 0<\gamma\le 1.
\]
其中 $\gamma$ 控制未来 reward 的权重。本节课的直觉讲解主要使用未折扣的总 reward。

\subsection{Step 3：优化 policy}

第三步是找参数 $\theta$，让 actor 与环境互动时得到的 return 尽可能大。因为 actor 的动作会影响后续 state，所以 $\theta$ 不只影响当前输出，也影响整条 trajectory 的分布。目标可写成
\[
J(\theta)=\E_{\tau\sim p_\theta(\tau)}[R(\tau)],
\qquad
\theta^*=\argmax_\theta J(\theta).
\]
其中 $p_\theta(\tau)$ 表示由 policy $\pi_\theta$ 与 environment 共同诱导出的 trajectory 分布。

\framefig{0.92}{figures/fig12_trajectory_sequence.jpg}
{actor 与 environment 交替产生 state、action 和 reward，形成一条 trajectory。}
{00:25:00--00:25:15}

\framefig{0.92}{figures/fig13_maximize_expected_return.jpg}
{RL 优化目标：调整 policy network 的参数，使 trajectory 的 return 越大越好。}
{00:27:15--00:27:30}

\subsection{本章小结}

RL 的三步骤与一般机器学习一致：定义 policy network、定义 return 目标、优化参数。差别在于目标函数不是静态数据集上的 loss，而是 policy 与环境互动后得到的期望 return；policy 的参数会改变未来数据分布，这让优化更难。

% =====================================================================
\section{为什么 RL 优化困难}
% =====================================================================

如果 environment 和 reward 都是可微、确定且已知的神经网络，那么我们似乎可以直接对 $J(\theta)$ 做反向传播。但现实中的 RL 不满足这个理想条件。environment 常常是黑盒，例如游戏机、真实机器人、围棋对手或在线用户；reward 函数也未必可微；更麻烦的是，互动过程本身有随机性。

\framefig{0.92}{figures/fig14_black_box_randomness.jpg}
{environment 和 reward 通常是黑盒，且 trajectory 采样带有随机性；同样的 policy 不一定每次产生同样结果。}
{00:29:45--00:30:00}

这种黑盒随机性意味着，不能简单地把整个环境接到神经网络后面做端到端微分。我们只能让 actor 去试、收集 trajectory、观察 reward，再根据这些经验调整 policy。这也解释了为什么 RL 的训练通常不稳定、样本效率低，而且实验结果方差很大。

\subsection{与 GAN 的类比}

课程用 GAN 作类比：在 GAN 里，generator 产生样本，discriminator 给出评价；训练 generator 时，我们希望 discriminator 的输出更高。RL 中 actor 有点像 generator，environment 和 reward 有点像 discriminator；我们调 actor，让 reward 更高。

\framefig{0.92}{figures/fig15_gan_analogy_policy_optimization.jpg}
{RL 优化可类比 GAN：actor 像 generator，environment 与 reward 像评价系统；但 RL 的评价系统通常是不可微黑盒。}
{00:31:00--00:31:15}

关键差别在于，GAN 的 discriminator 通常也是神经网络，我们知道它内部结构，能把梯度传回 generator；RL 的 environment 不一定是神经网络，也不一定可微。即便在电子游戏中，环境内部也可能有随机数；在围棋中，对手如何回应更不受我们控制。因此 policy optimization 需要专门方法，policy gradient 就是其中最基本的一类。

\begin{warningbox}{RL 代码容易跑，稳定结果不容易}
    老师提到 RL 常让人感觉又难又容易：网上有很多可运行代码，但同样 network 每次测试结果可能都不同。随机性和黑盒环境使 RL 实验方差很大，因此评估时通常需要多次 seed、平均回报和置信区间。
\end{warningbox}

\subsection{本章小结}

RL 优化困难的根源是：目标依赖整条 trajectory，trajectory 又由 policy 与黑盒环境共同产生；环境不可微、reward 延迟且采样有随机性。policy gradient 的目的，就是在这种不能直接反向传播的情况下，仍然找到调整 actor 的方向。

% =====================================================================
\section{Policy Gradient 的直觉：控制 actor}
% =====================================================================

课程最后进入 policy gradient，但没有从完整推导开始，而是用一个更贴近监督学习的角度讲：如果我们知道在某个 observation $s$ 下希望 actor 采取动作 $\hat a$，那就可以像训练分类器一样，用 cross-entropy 让 actor 更倾向于输出 $\hat a$。

\framefig{0.92}{figures/fig16_policy_gradient_section.jpg}
{进入 policy gradient：本节用“如何控制 actor”的角度解释 policy optimization。}
{00:35:45--00:36:00}

\subsection{让 actor 更常采取某个动作}

若希望 actor 在 state $s$ 下采取动作 $\hat a$，可以把 $\hat a$ 当成 label，计算 cross-entropy：
\[
e(s,\hat a;\theta)
=-\log \pi_\theta(\hat a\mid s).
\]
最小化 $e$ 会增大 $\pi_\theta(\hat a\mid s)$，也就是让 actor 更常采取 $\hat a$。

\framefig{0.92}{figures/fig17_cross_entropy_take_action.jpg}
{若希望 actor 在 observation $s$ 下采取动作 $\hat a$，可把动作当成 label，用 cross-entropy 训练。}
{00:37:00--00:37:15}

这和训练图像分类器很像：输入是图片，label 是类别；在 RL 中，输入是 observation，label 是希望采取的 action。区别是，RL 中这个 label 不是人类预先标好的，而是需要从互动经验中判断出来。

\subsection{让 actor 避免某个动作}

如果希望 actor 在某个 observation $s'$ 下不要采取动作 $\hat a'$，可以用相反方向的目标。直觉上，最小化
\[
-e(s',\hat a';\theta)
=\log \pi_\theta(\hat a'\mid s')
\]
会压低 $\pi_\theta(\hat a'\mid s')$，让 actor 不倾向于该动作。把“希望做”和“不希望做”放在一起，可以写成
\[
L(\theta)=e(s,\hat a;\theta)-e(s',\hat a';\theta).
\]

\framefig{0.92}{figures/fig18_take_and_avoid_actions.jpg}
{同时控制正向与负向行为：对希望采取的动作降低 cross-entropy，对不希望采取的动作提高 cross-entropy。}
{00:39:00--00:39:15}

\begin{warningbox}{这是 policy gradient 的直觉版}
    直接最小化负 cross-entropy 只是帮助理解“鼓励或抑制某个动作”的方向。实际 policy gradient 会从期望 return 的梯度推导出加权 log-probability 目标，并会加入 baseline、advantage、entropy regularization 等稳定训练技巧。
\end{warningbox}

\subsection{把行为资料整理成训练集}

如果我们收集到许多 $(s_n,a_n)$，并知道每个动作是“希望执行”还是“不希望执行”，就可以像分类任务一样训练 actor。课程用 $+1$ 表示希望做，用 $-1$ 表示不希望做：
\[
\mathcal{D}
=\{(s_n,a_n,y_n)\}_{n=1}^{N},
\qquad
y_n\in\{+1,-1\}.
\]
对应的直觉 loss 可以写成
\[
L(\theta)=\sum_{n=1}^{N}y_n\,e(s_n,a_n;\theta).
\]
若 $y_n=+1$，最小化会鼓励动作；若 $y_n=-1$，最小化会抑制动作。

\framefig{0.92}{figures/fig19_binary_training_data_for_actor.jpg}
{把 actor 控制问题整理成训练资料：某些 state-action pair 标成 Yes，某些标成 No。}
{00:41:45--00:42:00}

\subsection{从二元标签到 advantage 权重}

实际情况不只是“好”或“不好”。某个动作可能非常好、稍微好、稍微坏或非常坏。因此课程把 $+1/-1$ 推广成一个实数权重 $A_n$。$A_n$ 越大，越希望 actor 在 $s_n$ 下采取 $a_n$；$A_n$ 越小甚至为负，越希望 actor 避免该动作。

\framefig{0.92}{figures/fig20_advantage_weighted_training_data.jpg}
{每个 state-action pair 不只有二元好坏，而可以有强弱不同的权重 $A_n$。}
{00:42:30--00:42:45}

于是 actor 的训练目标可以写成
\[
L(\theta)=\sum_{n=1}^{N}A_n\,e(s_n,a_n;\theta)
=-\sum_{n=1}^{N}A_n\log \pi_\theta(a_n\mid s_n).
\]
最小化这个 loss 等价于：对 $A_n>0$ 的动作，提高其概率；对 $A_n<0$ 的动作，降低其概率。这个形式已经很接近 policy gradient 中常见的 advantage-weighted log-probability objective。

\framefig{0.92}{figures/fig21_weighted_loss_for_actor.jpg}
{用 $A_n$ 加权每个 cross-entropy 项：正权重鼓励动作，负权重抑制动作。}
{00:44:00--00:44:15}

\begin{importantbox}{policy gradient 的一行直觉}
    让 actor 回顾自己做过的动作：如果某个动作后来带来好结果，就增加它在相似状态下被采样的概率；如果带来坏结果，就降低它的概率。数学上，这会变成用回报或 advantage 加权的 $\log\pi_\theta(a\mid s)$。
\end{importantbox}

\subsection{这一讲停在哪里}

课程最后留下两个问题。第一，$A_n$ 要怎么产生？也就是如何判断某个 state-action pair 到底该被鼓励还是抑制，以及鼓励程度是多少。第二，训练资料 $(s_n,a_n)$ 从哪里来？它们显然不能像监督学习那样由人工完整标注，而要来自 actor 与环境的互动轨迹。

\framefig{0.92}{figures/fig22_unknown_source_of_training_data.jpg}
{本节结尾的问题：训练 actor 所需的 $(s_n,a_n)$ 与权重 $A_n$ 从哪里来？这正是后续 policy gradient 要解决的核心。}
{00:44:45--00:45:00}

\subsection{本章小结}

policy gradient 可以先用监督学习直觉理解：如果知道某个动作该做，就用 cross-entropy 提高它的概率；如果知道不该做，就降低它的概率；如果知道好坏程度，就用权重 $A_n$ 调整力度。真正的 RL 问题是，$A_n$ 和训练样本都必须从互动经验中估计，而不是由静态标签直接给出。

% =====================================================================
\section{总结与延伸}
% =====================================================================

本节课完成了强化学习入门的第一层框架。RL 仍然是机器学习三步骤：定义 actor，定义目标，优化参数。actor 是 policy network，输入 observation，输出 action 分布；environment 根据 action 变化，并回传 reward。一次完整互动形成 trajectory 或 episode，reward 总和称为 return。优化目标是最大化 policy 诱导出的期望 return：
\[
J(\theta)=\E_{\tau\sim p_\theta(\tau)}[R(\tau)].
\]

RL 与监督学习最大的不同，在于没有每一步的正确答案。label 隐藏在长期互动结果里，而且 action 会改变之后看到的数据。Space Invader 中 reward 可以相对频繁地出现；围棋中 reward 可能只在终局出现。这样的 delayed reward 让 credit assignment 变得困难，也让 policy optimization 不能简单套用普通 backpropagation。

从 policy gradient 的角度看，训练 actor 可以被理解为加权分类：对某个 state-action pair，如果权重 $A_n$ 为正，就提高该动作概率；如果 $A_n$ 为负，就降低该动作概率。于是
\[
L(\theta)=-\sum_{n=1}^{N}A_n\log\pi_\theta(a_n\mid s_n)
\]
成为一个很有解释力的形式。它把 RL 的“长期好坏”转换成对局部动作概率的鼓励或惩罚。

\subsection{下一步要解决的问题}

这节课刻意停在最关键的缺口上：我们还不知道 $A_n$ 怎么来，也不知道哪些 $(s_n,a_n)$ 应该进入训练集。后续 policy gradient 会说明，trajectory 是由当前 policy 采样出来的，而 $A_n$ 可以由 return、baseline 或 advantage function 估计。这会把“鼓励好动作、抑制坏动作”的直觉，变成可以实际训练的算法。

\begin{importantbox}{这一讲的最重要收获}
    不要把强化学习想成完全脱离前面课程的新世界。它仍然是在训练函数，只是训练信号来自与环境互动后的 reward。理解 actor、trajectory、return 和加权 cross-entropy，就抓住了 policy gradient 的入口。
\end{importantbox}

\end{document}
